% ─────────────────────────────────────────────────────────────────────────
% §5 Discussion
% ─────────────────────────────────────────────────────────────────────────
\section{Discussion}\label{sec:discussion}

The seven findings in §\ref{sec:results} resist summarization under a
single-axis alignment-culture theory. This section develops three
corresponding departures from standard framings in the public discourse
on LLM refusal behavior.

\subsection{Alignment Culture Is Not a One-Dimensional East-West Axis}\label{sec:discussion-axis}

A frequent simplification in popular and industry commentary on LLM
alignment treats national origin as the primary organizing axis:
\emph{Western} vendors are imagined to share an RLHF-driven
``liberal-democratic'' alignment posture and \emph{Chinese} vendors an
``authoritarian'' one. Under that framing, the observable behaviors of
\vendor{OpenAI}, \vendor{Gemini}, \vendor{DeepSeek}, and \vendor{Kimi}
should partition 2-2 along corporate geography.

Finding 1 refutes this prediction. The two Chinese-owned vendors produce
the single most distant pair of refusal distributions in the matrix
(\vendor{DeepSeek}-\vendor{Kimi} JSD 0.200, CI \ci{0.149}{0.256}). The
two Western-Chinese cross-pairs that the East-West framing predicts to
be distant, \vendor{OpenAI} vs \vendor{DeepSeek} and \vendor{Gemini} vs
\vendor{DeepSeek}, are in fact the closest non-Western pairings, with
JSDs an order of magnitude below the \vendor{DeepSeek}-\vendor{Kimi}
gap. Neither a geographic nor a regulatory-regime axis explains the
clustering that actually obtains.

Under a slightly more mechanistic view this is not surprising.
\vendor{DeepSeek}'s training recipe is publicly documented and
explicitly acknowledges RLHF-style post-training on curated preference
data in the tradition of InstructGPT and its successors. At the level
of distribution shape (the relative mix of hard, soft, and on-task
responses), \vendor{DeepSeek} tracking \vendor{OpenAI} and
\vendor{Gemini} is precisely what one would expect from a vendor that
has adopted contemporary RLHF conventions, regardless of regulatory
geography. \vendor{Kimi}'s departure is likewise not evidence of
``Chinese alignment'' in some monolithic sense; it is evidence of a
vendor-specific architectural choice to split the refusal workload
between an infrastructure-level filter and a permissive underlying
model (§\ref{sec:discussion-two-layer}). What Findings 1 and 5 jointly
refute is the framing itself, not any claim about a particular vendor's
alignment politics.

The closest comparable audit we are aware of, \citet{ko2026bilingual},
probes 17 LLMs on Taiwan sovereignty questions using parallel
Chinese/English queries and reports significant language-dependent
bias in 15 of 17. That study's finding, that the same vendor can
behave materially differently on sovereignty-adjacent content depending
on query language, is orthogonal to but compatible with our per-topic
and per-layer findings. Both studies suggest that sovereignty-topic
LLM behavior is richer than aggregate vendor characterizations allow.

\subsection{Topic \texorpdfstring{$\times$}{x} Vendor \texorpdfstring{$\times$}{x} Layer: A Three-Way Interaction}\label{sec:discussion-three-way}

The cleaner organizing principle that emerges from our data is a
three-way interaction between prompt \emph{topic}, vendor \emph{identity},
and the \emph{layer} at which refusal is expressed. The
aggregate-distribution similarity of §\ref{sec:results-finding1} and the
topic-stratified divergence of §\ref{sec:results-finding5} are not in
tension: they describe the refusal surface from two different
cross-sections. At the aggregate marginal of prompt topics, DeepSeek
shapes up like OpenAI/Gemini; conditional on the sovereignty topic,
DeepSeek is categorically more restrictive than any other vendor in
our panel.

This interactive structure implies that
\emph{topic-aggregated refusal metrics materially misrepresent vendor
behavior.} A downstream researcher using DeepSeek for non-sovereignty
Taiwan content (e.g., ethnic relations, historical events, policy
debates) would observe behavior essentially indistinguishable from
Western vendors at refusal-rate metrics commonly reported in the
alignment literature. The same researcher moving to sovereignty-adjacent content, which is
relevant for any simulation of Taiwanese political engagement with
cross-strait questions, would encounter a refusal regime that permits
on-task engagement 10\,\% of the time. That is a 40-percentage-point
shift invisible to aggregate reporting.

The layer dimension adds a second implication. \vendor{Kimi}'s aggregate
on-task rate of 83\% on labeled responses is the highest in our panel.
A researcher reading only the aggregate number and comparing it to the
public characterization of Kimi as a Chinese-vendor aligned product
might be surprised. The resolution is that the aggregate rate is
computed over prompts that \emph{reached} the language model ---
i.e., after the infrastructure filter removed 7\% of calls. The full
denominator, including filter rejections, is different in kind:
infrastructure rejections are visible to the caller as terminal errors
with filter-reason messages, not as outputs requiring content analysis.
Treating the two as commensurable categories within a single refusal
metric is a conceptual error that inflates Kimi's apparent
permissiveness.

Concretely, for vendor audit practice, the implications are:
\begin{itemize}
  \item \textbf{Report Layer 1 and Layer 2 refusal separately.} For a
    vendor with non-trivial Layer 1 filtering (only \vendor{Kimi} in our
    dataset, but the set may grow), the aggregate refusal rate conflates
    two phenomena with different auditability properties.
  \item \textbf{Stratify by topic.} Aggregate refusal rates can be
    near-identical for vendors whose topic-conditional behaviors differ
    by factors of 4--5 (DeepSeek vs Western on sovereignty). For any
    simulation study whose prompt distribution is not uniform over
    topics, aggregate rates are inadequate.
  \item \textbf{Stratify by prompt design intent}, not only by topic.
    The 8.0\% OT-expected block rate for Kimi is invisible in
    topic-level reporting (it contributes diffusely to the factual
    topic's 12.5\% rate) but is the most interpretable evidence for the
    \emph{Taiwan-statehood blocking} characterization of the filter.
\end{itemize}

\subsection{The Two-Layer Architecture as a Distinct Alignment Strategy}\label{sec:discussion-two-layer}

The clearest architectural contrast in our dataset is between
\vendor{DeepSeek}'s topic-specific RLHF strengthening on sovereignty
(sovereignty on-task 10\%, non-sovereignty 54\%) and \vendor{Kimi}'s
infrastructure-filter-plus-permissive-model combination
(sovereignty 83\% of labeled responses on-task, but with 7\% of all
sovereignty responses blocked upstream). These are two distinct engineering solutions to the same
alignment-policy objective (suppressing responses to
sovereignty-sensitive prompts that an alignment team would not endorse),
with meaningfully different properties for downstream auditability.

The DeepSeek-style strategy concentrates the refusal workload in a
single layer (the model's RLHF-shaped output distribution) and
consequently produces refusals that are (i) probabilistic (the same
prompt may on different samples yield a refusal or an engagement),
(ii) content-sensitive (prompt rephrasing is effective at influencing
the response), and (iii) often visually similar to engagement (soft
refusal, not hard refusal, dominates). The resulting auditor
must work at the content-analysis level, with all the rater-reliability
challenges documented in §\ref{sec:method-reliability}.

The Kimi-style strategy splits the refusal workload between a
deterministic Layer 1 and a permissive Layer 2. Layer 1 refusals are
visible to every API caller as explicit filter-reason errors; they
require no content analysis to identify and no inter-rater protocol to
measure. The trade-off is that Layer 1 filtering is maximally visible:
a downstream user who prompts \vendor{Kimi} with ``How many seats are
in Taiwan's legislature?'' (OT03) receives a filter error and can
immediately recognize it as a policy decision by the vendor, rather
than a hedged engagement that might be mistaken for a capability
limit.

Neither strategy is clearly preferable from a pure auditability
standpoint. The \vendor{Kimi} strategy is more transparent about its
hard refusals but can create user-perceived gaps on factual content;
the \vendor{DeepSeek} strategy preserves user-perceived coverage but
hides its refusals inside content-level reasoning. What our data
supports is the empirical observation that both strategies exist in
production at commercial scale and produce qualitatively different
audit signals, which should not be aggregated into a single ``refusal
rate''.

\subsection{The DeepSeek Paradox and Its Possible Origins}\label{sec:discussion-deepseek}

The single strongest signal in our dataset is the conjunction of
\vendor{DeepSeek}'s Western-shaped aggregate refusal distribution with
its sovereignty-specific 10.3\% on-task rate. We are not positioned from
this audit alone to adjudicate its mechanistic origin, but three
hypotheses are worth registering for future investigation.

First, DeepSeek's RLHF reward model may be directly shaped by a
topic-specific penalty for sovereignty-engaging outputs that does not
fire at comparable intensity on Western RLHF models. This would be
consistent with a training-data curation decision taken at the DeepSeek
organization and would predict that smaller or fine-tuning-modifiable
versions of the model exhibit reduced topic specificity.

Second, DeepSeek's pre-training data may under-represent a particular
class of content (e.g., Traditional Chinese political writing authored
from a Taiwan perspective) such that the model's sovereignty-topic
responses are both more error-prone and more defensively hedged, not
because the alignment layer targets the topic but because the base
model has lower competence on it. This would predict that
sovereignty-topic soft refusals exhibit semantic patterns consistent
with uncertainty-driven hedging rather than policy-driven refusal.

Third, the observation might be driven in part by a content-filter
analog operating at the \vendor{DeepSeek} API layer that we do not
detect because it manifests as soft-refusal output rather than
terminal filter errors. That is, \vendor{Kimi}'s filter and
\vendor{DeepSeek}'s filter may be functionally similar but surface
differently: \vendor{Kimi}'s appears as explicit API errors,
\vendor{DeepSeek}'s as RLHF-shaped hedged output. This would predict that on sovereignty-adjacent prompts,
DeepSeek's non-response content exhibits statistical regularities
suggestive of a deterministic filter gate upstream of the generation.
\citet{naseh2025r1dacted} report evidence consistent with this
hypothesis for \vendor{DeepSeek}'s R1 reasoning model (which we do not
audit here), characterizing it as the first instance of local censorship
alignment at the model layer that persists under private deployment.
Whether the same mechanism is operative in the non-reasoning
\texttt{deepseek-chat V3.2} endpoint we audit is an open question
that our instrumentation cannot answer.

We lack instrumentation to discriminate between these hypotheses from
the present dataset. All three would predict a meaningful
DeepSeek--sovereignty interaction; the present work's contribution is
to document that interaction rigorously and to point subsequent work
toward mechanistic characterization.

\subsection{Robustness to Model-Tier Asymmetry}\label{sec:discussion-robustness}

Section~\ref{sec:method-vendors} disclosed that our five-vendor panel
is not capability-matched: \vendor{OpenAI} and \vendor{Gemini} are
represented by mini/lite-class models, whereas \vendor{DeepSeek} and
\vendor{Kimi} are flagship serving endpoints. A natural alternative
explanation for Findings~1, 3, 5, and 7 is that model scale rather than
alignment policy drives the observed inter-vendor divergence. To test
this, we re-ran a stratified 40-prompt subset of the main experiment
under a flagship-tier configuration: \vendor{OpenAI} queried at
\texttt{gpt-4o}, \vendor{Gemini} at \texttt{gemini-2.5-flash}, and
\vendor{Grok} at \texttt{grok-3}. (\vendor{Gemini-2.5-pro} was
considered but requires reasoning mode, which would break the
non-reasoning consistency of our vendor configuration; we use
\texttt{gemini-2.5-flash}, the flagship non-reasoning variant in the
flash family.) \vendor{DeepSeek} and \vendor{Kimi} already used their
flagship endpoints in the main experiment and are unchanged. Paired
bootstrap CIs are computed on the subset using the same protocol as
§\ref{sec:method-stats}.

Figure~\ref{fig:tier-comparison} reports the comparison on the four
tier-sensitive findings. The n=40 subset produces wide confidence
intervals on any single-vendor statistic; we therefore focus on
qualitative direction and rank-ordering rather than on point-estimate
equality.

\begin{figure}[htbp]
  \centering
  \includegraphics[width=\textwidth]{fig7_tier_comparison.pdf}
  \caption{Production-tier (gray) vs flagship-tier (orange) comparison
  on four tier-sensitive headline findings. Error bars are 95\% paired
  bootstrap CIs (5{,}000 resamples, prompt-level pairing). Highlighted
  cells mark the critical \vendor{DeepSeek} × sovereignty cell (B) and
  the \vendor{DeepSeek}$\leftrightarrow$\vendor{Kimi} JSD pair (D) that
  anchor Findings 5 and 1 respectively.}
  \label{fig:tier-comparison}
\end{figure}

\paragraph{Finding~1 (refusal-distribution clustering) is robust.}
In the production tier restricted to the subset, the pairwise JSDs
preserve the main-experiment ordering:
\vendor{DeepSeek}$\leftrightarrow$\vendor{Kimi} is the largest
off-diagonal entry (0.209, 95\% CI \ci{0.097}{0.355}), while
\vendor{DeepSeek}$\leftrightarrow$\vendor{OpenAI} (0.074) and
\vendor{DeepSeek}$\leftrightarrow$\vendor{Gemini} (0.085) are an order
of magnitude smaller. In the flagship tier, the same ranking obtains:
\vendor{DeepSeek}$\leftrightarrow$\vendor{Kimi} remains the
maximum (0.236, CI \ci{0.115}{0.384}), with
\vendor{DeepSeek}$\leftrightarrow$\vendor{Gemini} actually
\emph{decreasing} (0.047, CI \ci{0.007}{0.098}) when \vendor{Gemini} is
queried at its flagship model. The shape-level clustering that refutes
the East--West dichotomy is therefore not produced by the mini/flagship
capability gap; it persists when the gap is eliminated.

\paragraph{Finding~5 (\vendor{DeepSeek} sovereignty collapse) is robust.}
In both tiers, \vendor{DeepSeek}'s sovereignty on-task rate on the
7-prompt sovereignty subset of the 40-prompt subset is
\textbf{0.0\%} (0/7 in each tier). No other vendor's sovereignty
on-task rate drops below 28\% in either tier. Because the
\vendor{DeepSeek} model identifier is identical across tiers, this
cross-tier stability is trivial for \vendor{DeepSeek} specifically;
what the flagship subset adds is evidence that none of the three
\emph{upgraded} vendors (\vendor{OpenAI}, \vendor{Gemini},
\vendor{Grok}) produces a comparable sovereignty collapse when
queried at flagship scale. Flagship \vendor{Gemini}'s sovereignty
on-task rate drops from 71\% to 29\% in our subset (small-$n$ noise
or a real shift we cannot resolve at this sample size), but even at
29\% it is an order of magnitude above \vendor{DeepSeek}'s 0\%.
The sovereignty-specific dampening of \vendor{DeepSeek} is not an
artifact of querying it against smaller Western models.

\paragraph{Finding~3 (\vendor{Grok}/\vendor{Kimi} low-refusal regime)
is robust.}
Aggregate refusal rates in the subset preserve the production-tier
ordering: \vendor{DeepSeek} highest (62.5\% in both tiers),
\vendor{Kimi} lowest (12.5\% prod, 10.0\% flagship), \vendor{Grok}
second-lowest (17.5\% prod, 22.5\% flagship). The
\vendor{Grok-3}-vs-\vendor{grok-4-fast-non-reasoning} comparison is
particularly informative: despite the model-scale difference, Grok
produces qualitatively similar balanced-hedge responses on
partisan-political prompts (observed in side-by-side inspection during
labeling). The low-refusal regime is a vendor-alignment feature,
not a model-scale feature.

\paragraph{Finding~7 (two-tier HR$\to$SR elasticity) is robust.}
\vendor{DeepSeek}'s elasticity is literally identical (9.8\,pp in both
tiers, same model). \vendor{OpenAI}/\vendor{Gemini}/\vendor{Grok}
flagship elasticities fall in 26.8\,pp--42.9\,pp, overlapping the main
experiment's responsive-tier range. CIs are wide in the 40-prompt
subset and many pairwise differences are not statistically distinct,
but no flagship vendor collapses to a low-Δ regime that would
challenge the responsive/non-responsive two-tier picture.

\paragraph{Finding~2 is not tested here.}
The \vendor{Kimi} Taiwan-statehood filter (§\ref{sec:results-finding2})
is vendor-infrastructure-level and model-independent;
\vendor{Kimi}'s model identifier is identical between main and
flagship tiers, so there is no tier contrast to test. The 0/40
API-blocks observed in the flagship subset reflect Poisson variation
at n=40 with a 7\% main-experiment filter rate (expected
$\approx 2.8$, $P(X=0 \mid \lambda=2.8) \approx 6\%$) and does not
constitute evidence for or against the main Finding~2 characterization.

\paragraph{Overall.}
The size-confound alternative explanation for the paper's headline
findings is not supported by the flagship-tier sensitivity subset.
The four tier-sensitive findings retain their qualitative character
when \vendor{OpenAI}, \vendor{Gemini}, and \vendor{Grok} are queried
at capability-matched flagship endpoints. We report this as
corroborating robustness rather than as confirmation; the 40-prompt
subset is modest, and a scaled-up replication at
$N \geq 200$-prompt flagship fetch would provide tighter CIs on each
finding. All data, labels, and comparison JSON from this sensitivity
subset are released alongside the main dataset.

\subsection{Implications for LLM Agent Simulation Studies}\label{sec:discussion-simulation}

This audit was originally motivated by a larger project using LLM
agents to simulate Taiwan voter populations. The findings are
directly consequential for that class of application. A
simulation study using a single vendor for Taiwan-political agents
will produce results that are vendor-specific in ways not always
acknowledged in the computational social science literature. Two
concrete risks deserve explicit highlighting:

\begin{itemize}
  \item A simulation restricted to \vendor{DeepSeek} would produce
    populations whose sovereignty-topic discourse is dominated by
    soft refusal or on-task engagement shaped by the model's
    10\%-sovereignty-on-task regime. Results on cross-strait attitudes
    would systematically under-represent engagement with the topic,
    and soft refusal content may be misinterpreted as agent-level
    indecision rather than vendor-level alignment shaping.
  \item A simulation restricted to \vendor{Kimi} would silently drop
    7\% of sovereignty-adjacent prompts (infrastructure-level), with
    the remaining responses skewing sharply toward on-task engagement
    at the post-filter language model. Aggregated across an agent
    day, this could manifest as vendor-specific sparsity in
    sovereignty discussion combined with (selectively) higher
    compliance on what remains.
\end{itemize}

Our recommended mitigation is simple: report vendor as a first-class
variable in any Taiwan-political agent simulation, and where feasible
run the same simulation over multiple vendors to quantify
vendor-dependent effects. The audit data in this paper provides the
prior distribution under which such multi-vendor comparisons should
be interpreted.
